\documentclass[11pt,a4paper]{article}

% ── Packages ──────────────────────────────────────────────────────────────────
\usepackage[a4paper, top=18mm, bottom=18mm, left=22mm, right=22mm]{geometry}
\usepackage{times}
\usepackage{microtype}
\usepackage{booktabs}
\usepackage{tabularx}
\usepackage{makecell}
\usepackage{ragged2e}
\usepackage{array}
\usepackage{multirow}
\usepackage{enumitem}
\usepackage[table]{xcolor}
\usepackage{titlesec}
\usepackage{fancyhdr}
\usepackage{hyperref}
\usepackage{amsmath}
\usepackage{float}
\usepackage{caption}
\usepackage{parskip}
\usepackage{longtable}

% ── Colours ───────────────────────────────────────────────────────────────────
\definecolor{iitblue}{RGB}{26,35,126}
\definecolor{goodgreen}{RGB}{27,94,32}
\definecolor{warnred}{RGB}{183,28,28}
\definecolor{rowgray}{RGB}{245,245,255}

% ── Section formatting ────────────────────────────────────────────────────────
\titleformat{\section}
  {\normalfont\large\bfseries\color{iitblue}}{}{0em}{}[{\color{iitblue}\titlerule[0.7pt]}]
\titleformat{\subsection}
  {\normalfont\normalsize\bfseries\color{iitblue}}{}{0em}{}
\titlespacing*{\section}{0pt}{10pt}{5pt}
\titlespacing*{\subsection}{0pt}{7pt}{3pt}

% ── Header / Footer ───────────────────────────────────────────────────────────
\pagestyle{fancy}
\fancyhf{}
\renewcommand{\headrulewidth}{0.4pt}
\setlength{\headwidth}{\textwidth}
\fancyhead[L]{\small\color{iitblue}\textbf{Midterm Report:} DP LLM Fine-Tuning}
\fancyhead[R]{\small IIT Guwahati}
\fancyfoot[C]{\small\thepage}

% ── List spacing ──────────────────────────────────────────────────────────────
\setlist[itemize]{itemsep=2pt, topsep=3pt, leftmargin=16pt}
\setlist[enumerate]{itemsep=2pt, topsep=3pt, leftmargin=18pt}

% ── Table helpers ─────────────────────────────────────────────────────────────
\newcolumntype{C}[1]{>{\centering\arraybackslash}p{#1}}
\newcolumntype{L}[1]{>{\raggedright\arraybackslash}p{#1}}
\renewcommand{\arraystretch}{1.2}
\captionsetup{font=small, labelfont=bf, skip=4pt}

\setlength{\parskip}{4pt}
\setlength{\emergencystretch}{2em}
\setlength{\headheight}{14pt}
\setlength{\parindent}{0pt}

% ══════════════════════════════════════════════════════════════════════════════
\begin{document}

% ── Title block ───────────────────────────────────────────────────────────────
\begin{center}
  {\LARGE\bfseries\color{iitblue}
    Privacy-Utility Trade-off in\\[3pt]
    Differentially Private LLM Fine-Tuning}\\[6pt]
  {\large\textbf{Midterm Progress Report}}\\[4pt]
  {\normalsize Hardware aware DL (Sem VII) $\;\cdot\;$ IIT Guwahati $\;\cdot\;$ October 2026}\\[4pt]
  {\normalsize
    G Jayagiri $\;\cdot\;$
    Vasudev Krishna $\;\cdot\;$
    Uditya Shekawat $\;\cdot\;$
    Arkaprava Paul $\;\cdot\;$
    Harshit Goel}
\end{center}
\vspace{1mm}
{\color{iitblue}\hrule height 1.5pt}
\vspace{4mm}

% ══════════════════════════════════════════════════════════════════════════════
\section{Project Overview \& Foundation Papers}

\subsection{What This Project Is About}

Large language models (LLMs) are increasingly fine-tuned on private and sensitive data ---
patient records, legal documents, personal communications --- to adapt them to specialized
downstream tasks. However, fine-tuning creates a privacy risk: the resulting model can
\emph{leak} information about individual training examples through \textbf{Membership
Inference Attacks (MIA)}, where an adversary queries the model to determine whether a
specific sample was part of its training set.

\textbf{Differential Privacy (DP)} provides a mathematical guarantee against such leakage
by injecting calibrated Gaussian noise into the gradient update at each training step
(DP-SGD). The formal guarantee is parameterized by $\varepsilon$ (privacy budget) and
$\delta$ (failure probability): smaller $\varepsilon$ means stronger privacy. The
fundamental problem, however, is that DP-SGD significantly \emph{degrades model accuracy}
because the noise overwhelms the gradient signal, especially for large models with many
trainable parameters.

This project asks a precise mechanistic question:

\begin{center}
\emph{``Does the location of the Transformer layers we choose to fine-tune --- early,
middle, or late in the network --- affect how well a model learns a task under DP,
while maintaining the same formal privacy guarantee?''}
\end{center}

The hypothesis is that \textbf{late layers} (closer to the output) encode more
task-specific, class-discriminative representations, and therefore provide a stronger
gradient signal per trainable parameter under DP noise. By concentrating updates on
these layers, we expect to achieve better accuracy with the same privacy budget compared
to fine-tuning early layers or the entire network. We also study whether a
\emph{DP-aware automatic layer selection algorithm} can match or exceed the best
hand-chosen layer set.

The project additionally evaluates whether the choice of fine-tuned layers affects
\emph{empirical} privacy leakage --- i.e., whether an attacker using the strongest
known MIA can infer more about training data from a model trained on late layers vs.\
early layers.

\subsection{Foundation Papers and How They Are Integrated}

\textbf{Paper 1 --- DP-SelFT (Sha et al., arXiv:2605.17432, 2025).}
This paper proposes a framework for \emph{Differentially Private Selective Fine-Tuning}
(DP-SelFT). The key insight is that not all layers are equally useful or equally
vulnerable to DP noise. DP-SelFT selects which layers to train using a three-stage
pipeline: (1) generate a synthetic dataset privately using a commercial LLM API with a
small privacy budget $\varepsilon_{\text{sel}}$; (2) temporarily train each candidate
layer subset on this synthetic data under \textbf{worst-case perturbations} of the same
scale as the downstream DP noise --- this favours layers that are robust to noisy
updates, not just layers that appear informative under clean training; (3) fine-tune only
the selected layers on private data using DP-AdamW. DP-SelFT achieves +2--3\% accuracy
over full DP fine-tuning at $\varepsilon = 5$ without any additional privacy cost for
selection.

\textbf{Paper 2 --- Marek et al. (arXiv:2606.09401, ICLR 2026).}
This paper benchmarks \emph{empirical privacy leakage} under DP LLM adaptations using
the strongest known attack: offline \textbf{Robust Membership Inference Attack (RMIA)}
with a shadow model trained on a disjoint split of the same data. The key finding is
that the formal $\varepsilon$ guarantee alone does not determine practical leakage ---
the \emph{distributional closeness} between pretraining data and fine-tuning data is the
dominant factor. At $\varepsilon = 8$, in-distribution data yields MIA AUC of 0.7--0.9;
out-of-distribution data is significantly safer. The paper also proposes a holistic
\emph{pretrain-adapt} privacy auditing framework covering four stages of the LLM
lifecycle.

\textbf{Integration.}
Our project sits at the intersection of both papers. Paper~1 addresses \emph{which
parameters to train} and how to select them under DP constraints; Paper~2 provides the
\emph{audit methodology} to measure whether those choices affect real-world leakage.
Specifically, we use Paper~2's \textbf{swap-pair shadow-model RMIA protocol} to audit
the configurations defined by Paper~1's layer-selection framework. We train two models,
$M_A$ on split $D_A$ and $M_B$ on split $D_B$, so that each model acts as the other's
shadow reference. This lets us measure per-layer-configuration empirical MIA AUC at the
same $\varepsilon$ budget --- a direct experiment that neither paper individually
conducts.

% ══════════════════════════════════════════════════════════════════════════════
\section{Achievements to Date}

\subsection{Infrastructure Built}

A complete end-to-end research pipeline was designed and implemented from scratch:
\begin{itemize}
  \item Training engine (\texttt{src/train.py}): step-based DP-SGD training loop using
        Opacus, with RDP accounting, \texttt{BatchMemoryManager} for memory-efficient
        per-sample gradients, gradient telemetry (clipping fraction, noise multiplier
        $\sigma$), and logit extraction for all audit splits after training.
  \item Model management (\texttt{src/models.py}): supports selective layer freezing
        (by index, position label, or ``head only''), embedding freezing, and
        trainable-parameter reporting.
  \item Layer selection (\texttt{src/select\_layers.py}): implements DP-SelFT Stages
        2--3 (DP histogram over a public candidate pool, worst-case perturbation ranking),
        outputting a \texttt{selection.json} with layer ranking and $\varepsilon_{\text{sel}}$.
  \item Audit engine (\texttt{src/audit.py}): implements loss-threshold baseline,
        continuous RMIA ratio attack, and offline Marek Alg.~1 (RMIA), reporting AUC,
        TPR@1\%FPR, TPR@5\%FPR, calibration gap, and generalization gap.
  \item Experiment orchestrator (\texttt{run\_experiment.py}): automated 8-configuration
        grid with the swap-pair protocol, CSV result aggregation, and Kaggle GPU runner.
\end{itemize}

\subsection{Experiments Completed}

Two full experiment runs were completed on Kaggle GPU (RoBERTa-base, SST-2, $N{=}1024$,
$\varepsilon{=}5$, $\delta{=}10^{-5}$). The first run (v1) revealed a critical training
failure in all DP configurations due to a noise multiplier of $\sigma{\approx}2.06$
overwhelming gradient signal. Hyperparameter fixes were applied and the full grid was
re-run (v2). Results are shown in Table~\ref{tab:results}.

\begin{table}[htbp]
\centering
\small
\caption{Validation accuracy and membership-inference AUC for v1 (failed) and v2 (fixed).}
\label{tab:results}
\setlength{\tabcolsep}{4pt}
\renewcommand{\arraystretch}{1.15}
\begin{tabularx}{\textwidth}{@{}>{\RaggedRight\arraybackslash}X
                                >{\centering\arraybackslash}p{1.35cm}
                                >{\centering\arraybackslash}p{1.35cm}
                                >{\centering\arraybackslash}p{1.25cm}
                                >{\centering\arraybackslash}p{1.25cm}
                                >{\centering\arraybackslash}p{1.15cm}@{}}
\toprule
\textbf{Configuration} & \textbf{Layers} & \textbf{Trainable params} &
\textbf{Acc.\ v1} & \textbf{Acc.\ v2} & \textbf{MIA AUC} \\
\midrule
Baseline clean (no DP) & All & 100\% & 89.8\% & 89.0\% & 0.67 \\
\rowcolor{rowgray}
Baseline full DP & All & 100\% & 50.0\% & 81.0\% & 0.51 \\
Head only DP & Head & 0.48\% & 50.0\% & 58.0\% & 0.52 \\
\rowcolor{rowgray}
Selective early $k=3$ & Layers 0--2 & 17.5\% & 50.0\% & 72.2\% & 0.52 \\
Selective mid $k=3$ & Layers 4--6 & 17.5\% & 50.0\% & 83.4\% & 0.50 \\
\rowcolor{rowgray}
Selective late $k=3$ & Layers 9--11 & 17.5\% & 50.0\% & 84.8\% & 0.51 \\
DP-SelFT $k=3$ & Selected (3, 6, 4) & 17.5\% & 50.0\% & 82.9\% & 0.51 \\
\rowcolor{rowgray}
Control $H'$ & Layers 9--11 & 17.5\% & 50.0\% & 84.7\% & 0.51 \\
\bottomrule
\end{tabularx}
\begin{minipage}{\textwidth}
\vspace{2pt}
\footnotesize \textit{Note:} The DP guarantee is maintained in v2 runs
($\varepsilon_{\mathrm{spent}}\leq 5.0$). AUC values near 0.5 indicate near-random
membership inference. The v1 results are included to document the initial training issue.
\end{minipage}
\end{table}

\subsection{Key Research Findings}

\begin{itemize}
  \item \textbf{Layer location matters, even at constant parameter count.} All selective
        k3 configs train exactly 17.5\% of parameters, yet accuracy ranges from 72.2\%
        (early) to 84.8\% (late) --- a 12.6-point gap attributable solely to layer
        location.
  \item \textbf{Late layers are most efficient under DP.} Layers 9--11 reach 84.8\%,
        only 4.2 points below the non-private ceiling (89.0\%), while training just
        17.5\% of parameters.
  \item \textbf{DP-SelFT validates the selection approach.} The automatic DP-aware
        selection (82.9\%) closely matches the best heuristic (late layers, 84.7\%) at
        a cost of only $\varepsilon_{\text{sel}}{=}0.3$ of the overall budget.
  \item \textbf{Formal DP is empirically confirmed.} All DP MIA AUC values are
        ${\approx}0.49$--$0.53$ (near-random), confirming the attacker gains no
        meaningful advantage at $\varepsilon{=}5$, regardless of layer choice.
  \item \textbf{Head-only DP is a negative result.} 592k parameters is insufficient
        capacity to learn a classification task under DP noise (58\% accuracy).
\end{itemize}

% ══════════════════════════════════════════════════════════════════════════════
\section{Pending Limitations \& Second-Half Roadmap}

\subsection{Critical Gaps (Affect Scientific Validity)}

\begin{table}[htbp]
\centering
\caption{Critical limitations and planned fixes}
\setlength{\tabcolsep}{5pt}
\begin{tabularx}{\textwidth}{@{}C{0.55cm} >{\RaggedRight\arraybackslash}X >{\RaggedRight\arraybackslash}X @{}}
\toprule
\textbf{\#} & \textbf{Limitation} & \textbf{Planned Fix} \\
\midrule
L1 & Single seed (42); SE $\approx\pm1.3\%$ on AUC; no confidence intervals &
     Run seeds 43, 44 $\to$ 3-seed mean $\pm$ SD \\
\rowcolor{rowgray}
L2 & No $\varepsilon$ sweep; Pareto curve is one point at $\varepsilon{=}5$ &
     Add $\varepsilon \!\in\! \{1, 8, \infty\}$; location effects become measurable at $\varepsilon{=}1$ \\
L3 & SST-2 phrase-family contamination: 83.4\% of items are sub-phrases; may artificially suppress MIA AUC &
     Group-aware one-item-per-family data splitting in \texttt{src/data.py} \\
\rowcolor{rowgray}
L4 & Head warmup uses private data ($D_A$/$D_B$) before PrivacyEngine; DP guarantee slightly overstated &
     Switch warmup to public proxy data (IMDb / Yelp) \\
\bottomrule
\end{tabularx}
\end{table}

\subsection{Significant Gaps (Affect Claim Strength)}

\begin{table}[htbp]
\centering
\caption{Gaps affecting strength of conclusions}
\setlength{\tabcolsep}{5pt}
\begin{tabularx}{\textwidth}{@{}C{0.55cm} >{\RaggedRight\arraybackslash}X >{\RaggedRight\arraybackslash}X @{}}
\toprule
\textbf{\#} & \textbf{Limitation} & \textbf{Planned Fix} \\
\midrule
L5 & $k$-layer fraction sweep not done; ``how many layers?'' unanswered &
     Configs for $k \!\in\! \{1,2,3,6,9,12\}$ at $\varepsilon{=}5$, heuristic and DP-SelFT \\
\rowcolor{rowgray}
L6 & No early stopping; training loss varies 0.03--2.43 across runs &
     Validation-loss patience-based stopping; per-step loss curves \\
L7 & Single dataset (SST-2); generalisability unknown &
     Add MNLI or QQP as secondary dataset check \\
\bottomrule
\end{tabularx}
\end{table}

\subsection{Timeline for Remaining Work}

\begin{table}[htbp]
\centering
\caption{Second-half schedule}
\setlength{\tabcolsep}{8pt}
\begin{tabularx}{\textwidth}{@{}C{1.05cm} >{\RaggedRight\arraybackslash}X >{\RaggedRight\arraybackslash}X @{}}
\toprule
\textbf{Week} & \textbf{Task} & \textbf{Expected Outcome} \\
\midrule
1--2 & Fix family-aware data splitting (L3); switch warmup to public data (L4) &
       Audit integrity assured; DP guarantee tight \\
\rowcolor{rowgray}
3--4 & $\varepsilon \!\in\! \{1,8,\infty\}$ sweep + seeds 43,\,44 (L1, L2) &
       Privacy-utility Pareto frontier + confidence intervals \\
5--6 & $k$-fraction sweep (L5); add early stopping (L6) &
       Utility vs.\ trainable-parameter curve \\
\rowcolor{rowgray}
7   & Secondary dataset experiment (L7); final result aggregation &
       Generalisability evidence \\
8   & Final report and presentation preparation &
       Complete project submission \\
\bottomrule
\end{tabularx}
\end{table}

\bigskip
\textbf{Summary.}
The core infrastructure, training pipeline, and initial results are complete. The primary
midterm finding --- that selective fine-tuning of late Transformer layers (9--11) under
DP-SGD achieves 84.8\% validation accuracy compared to the 89.0\% non-private ceiling,
while maintaining MIA AUC $\approx 0.51$ (near-random) --- is a substantive and
reproducible result. The second half of the project will harden these findings with
multi-seed reproducibility, an $\varepsilon$-sweep to map the full Pareto frontier, and
the $k$-fraction study needed to completely answer the research question.

\end{document}